\section{Fractional source covers and a paid deficit}\label{sec:source-covers}

We now improve the geometric term in \cref{thm:cyclic-freezing}.
The recovery bank remains the complete bank of \cref{thm:exact-source-bank}.
The covers introduced here are witnesses used to estimate its raw source
sets; they are not additional records available to the machine.

\subsection{Parameters and complete core covers}

Fix an admissible geometry, a selected $\mathrm Q_2$ subcore of at most $J$
families, and at most $q$ complete exceptional families. Choose integers
$K,R\ge0$, $J\ge\max(K,R)$ and $M\ge\max(2,R+1)$, distinguish at most $K$ selected
core families, protect at most $R$ selected core families, and choose a legal rank phase
$\rho$. All distinguished families are among those selected before the
phase; they need not be protected. Write
\begin{equation}\label{eq:cover-parameters}
 \ell=2M+3,\qquad C_0=16+24M(M-1),\qquad C_q=C_0(1+3C_0)^q.
\end{equation}
Let $A_\rho(u)$ be the complete orbit of the retained core at time $u$,
and $O_\rho(t)$ the complete orbit with the selected exceptions included.
The raw source set is $S_\rho(t)=x(O_\rho(t))$. Put
\begin{equation}\label{eq:alpha-tau}
 \alpha L=\sum_{G\text{ selected core, not distinguished}}
              \Var(x;[p_G,U_G]),\qquad
 \tau=\sum_{G\text{ selected core, not protected}} |G|.
\end{equation}
These sums use original full spans and hop counts, before rank deletion;
in particular they do not depend on $\rho$.

\begin{definition}[Fractional complete-core cover]\label{def:fractional-core-cover}
\formal{SourceCovers/CoreFractionalCover.lean}{SourceCovers.coreRawSets,SourceCovers.fractionalCoreCover,SourceCovers.fractionalCoreCover_exists_optimal,SourceCovers.one_le_fractionalCoreCover_toReal,SourceCovers.coreRawSets_phase_shadowOrbit_eq,SourceCovers.fractionalCoreCover_ne_top_iff_subset,SourceCovers.fractionalCoreCover_exactSourceBank_ne_top}
For a nonempty finite set $S_0\subseteq\{0,\ldots,L-1\}$, define
$\kappa_\rho(S_0)$ as the minimum of
\[
 \sum_{u=0}^{T-1}\lambda_u
 \quad\text{subject to}\quad
 \lambda_u\ge0,\qquad
 \sum_{u:\,a\in x(A_\rho(u))}\lambda_u\ge1\quad(a\in S_0).
\]
If no cover exists, set $\kappa_\rho(S_0)=+\infty$.
Covering orbits may be rooted at any computation time; their raw sets
may contain points outside $S_0$. A finite cover has an optimum, since
weights may be truncated at one and the resulting feasible set is
compact. It has weight at least one. In fact a finite cover exists
exactly when $S_0\subseteq\{x(u):0\le u<T\}$: every complete core
orbit contains its root and only earlier times. Thus unit weights at
all computation times suffice whenever the condition holds, including
for $S_0=S_\rho(t)$ at every computation time $t$.
\end{definition}

\begin{lemma}[Fractional interval covers]\label{lem:interval-rounding}
\formal{SourceCovers/IntervalCover.lean}{SourceCovers.exists_cover_packing,SourceCovers.exists_subcover_le_weight}
Let $S$ be a finite subset of a linearly ordered set, and let
$(A_i)_{i\in I}$ be a finite family of order-convex sets covering $S$.
Order-convex means that each set contains every point between two of
its points. There exist $C\subseteq I$ and $P\subseteq S$ such that
\[
 S\subseteq\bigcup_{i\in C}A_i,\qquad |C|=|P|,\qquad
 |P\cap A_i|\le1\quad(i\in I).
\]
Consequently, nonnegative weights $w_i$ satisfying
$\sum_{i:\,x\in A_i}w_i\ge1$ for every $x\in S$ admit an integral
subcover of cardinality at most $\sum_{i\in I}w_i$.
The target $S$ may be empty. In particular, this applies to closed,
open or half-open real intervals without a boundedness assumption.
\end{lemma}
\begin{outline}
Induct on the finite target. For a nonempty target let $x$ be its
least point, and let $y$ be the greatest target point sharing an
$A_i$ with $x$. Choose such a set $A_j$ containing $x,y$. Order
convexity makes $A_j$ cover all target points through $y$. No set
containing $x$ meets the residual target above $y$, by the choice of
$y$. Apply induction there after removing $j$ from the index family,
then add $j$ to the subcover and $x$ to the packing. Both cardinalities
increase by one, and every original set still contains at most one
packing point. The empty target uses two empty sets.

Fractional coverage implies ordinary coverage. Sum its inequalities
over the packing and interchange the finite sums. Each set's weight
occurs at most once, so $|C|=|P|\le\sum_iw_i$. The proof uses only
finite extrema, order convexity and finite sums; the shared library
supplies the declarations, which this paper re-exports.
\end{outline}
This rounding step preserves the total fractional weight in the
raw-source gap bound below. It does not alter the complete recovery bank.

\begin{lemma}[A fractional raw-source gap bound]\label{lem:fractional-gap}
\formal{SourceCovers/FractionalGap.lean}{SourceCovers.fractional_phase_gap,SourceCovers.fractional_core_gap,SourceCovers.fractional_core_gap_of_union_variation,SourceCovers.fractional_core_gap_of_weights}
With the parameters above, any nonempty $S_0$ with finite
$\kappa_\rho(S_0)$ satisfies
\begin{equation}\label{eq:fractional-gap}
 \gaparea(S_0)\ge
 \frac{(1-\alpha)_+^2}{2(2K+1)\kappa_\rho(S_0)},
 \qquad (z)_+=\max(z,0).
\end{equation}
Here $\gaparea(S)=\sum_g |g|^2/(2L^2)$, with the sum over the circular
gaps of the distinct points in $S$; a singleton has one gap of length $L$.

A sharper form retains actual component counts. Join consecutive
computation times across the elementary edges in the union $B$ of the
original nondistinguished full spans. Let $m_\rho(u)$ count the resulting
time components met by $A_\rho(u)$ and put
$V_B=\sum_{e\in B}|x(e+1)-x(e)|$. For any feasible cover weights,
\begin{equation}\label{eq:fractional-gap-incidence}
 \gaparea(S_0)\ge
 \frac{(1-V_B/L)_+^2}{2\sum_{u=0}^{T-1}m_\rho(u)\lambda_u}.
\end{equation}
Here $V_B\le\alpha L$ and $m_\rho(u)\le2K+1$.
The same conclusions hold for any retained subset of the selected
$\mathrm Q_2$ core. No rank phase or unit-speed assumption on $x$ is
needed for this estimate.
\end{lemma}
\begin{outline}
Label a time $t$ by the number of elementary edges before $t$ that are
outside $B$. Equal labels define contiguous time components, including
isolated times. Their internal edge sets partition $B$ within the
computation horizon. The interval
between the minimum and maximum raw coordinates on each component has
length at most its internal variation. Thus the sum of all component
interval lengths is at most $V_B$, with every elementary edge charged once.

A nondistinguished shadow stays in one component, since every edge
between its endpoints belongs to its original full span. A complete
core orbit has at most two initial successors and thereafter deterministic
continuation. Along either branch, a label change uses a distinguished
family whose owner precedes the current time. The new time is at most
that owner, so this family cannot pay for another change on the branch.
Counting the shared root label once gives $m_\rho(u)\le1+2K$.
The formal proof keeps the smaller count of retained distinguished
owners preceding the root before relaxing it to $K$.

Give each component interval the sum of the weights of the covering
orbits that meet it. These weights cover $S_0$ and have total
$\sum_um_\rho(u)\lambda_u$. By \cref{lem:interval-rounding}, they
admit an integral subcover with at most that many intervals. Their
total length is at most $V_B$.

Sort the distinct target points and append the first point plus $L$.
Call a consecutive gap exposed if no chosen interval contains both
endpoints. One interval can contain the left endpoint of at most one
exposed gap: if it contained two, it would contain both endpoints of
the earlier gap. Since every left endpoint is covered, the exposed
gap count is at most the subcover size. The other gaps have disjoint
elementary edge sets inside the covering intervals, so their total
length is at most $V_B$. The exposed gaps therefore have total length
at least $(L-V_B)_+$. Cauchy--Schwarz, followed by division by $2L^2$,
gives \cref{eq:fractional-gap-incidence}. The unrolled last gap handles
the circle seam and singleton target without an extra component.

Finally take an attained minimum cover, use $V_B\le\alpha L$ and
$m_\rho(u)\le2K+1$, and obtain \cref{eq:fractional-gap}. All extrema,
gap lengths and component incidences in this argument are constructed
from the actual time path and the complete core orbits.
\end{outline}
This supplies a verified geometric credit for the paid-deficit argument.
The shared-root count improves the former factor $4(K+1)$ to $2(2K+1)$;
the quantitative consequences below use the improved factor. It does not
supply economical covers for arbitrary native geometries.

\subsection{Anchored approximation}

We represent finite circular geometry by a common subdivision containing
all points and arc endpoints. Its elementary edges have real lengths
$w_e\ge0$ summing to $L>0$. A closed arc contains consecutive edges and
both endpoint vertices; its union length is the sum of the lengths of
the edges in that union. Thus an edge covered several times is counted
once. Subdivision indices are labels, not physical integer coordinates.
The formalization identifies each weighted gap with its real endpoint
displacement, including the gap across the chosen seam. Taking all edge
lengths equal to one recovers the raw-coordinate definition used above.
Write $\psi(e)=e-e^2/2$.

\begin{lemma}[Anchored erosion of circular gaps]\label{lem:anchored-erosion}
\formal{SourceCovers/AnchoredErosion.lean}{SourceCovers.penalty, SourceCovers.area_lower_bound, SourceCovers.area_lower_bound_linear, SourceCovers.anchored_displacement_area_lower_bound, SourceCovers.anchored_point_area_lower_bound}
Let $S_0,S$ be nonempty finite point sets on the circle of circumference
$L$, represented in a common subdivision as above. Suppose that
$S\subseteq S_0\cup\bigcup_{i=1}^m I_i$, where the $I_i$ are closed
circular arcs, each containing a point of $S_0$. If their union has
length at most $E$, set $e=\min\{E/L,1\}$. Then
\begin{equation}\label{eq:anchored-erosion}
 \gaparea(S)\ge\gaparea(S_0)-\psi(e)
             \ge\gaparea(S_0)-E/L.
\end{equation}
The statement permits points of $S_0$ outside $S$, overlapping arcs,
singleton arcs and full-circle arcs. For $0\le E\le L$, the loss bound is
$E/L-E^2/(2L^2)$.
\end{lemma}
\begin{outline}
Let $U$ be the union of covered elementary edges and let
$C=\sum_{a\in U}w_a$. Label each edge by the gap containing it, first
for $S_0$ and then for $S$. Two uncovered edges in the same $S_0$ gap
must share an $S$ gap. Indeed, any intervening $S$ point outside $S_0$
lies in a covering arc. That arc also contains an $S_0$ point outside
the intervening interval, so it must cross one of the two supposedly
uncovered edges. This argument uses actual arc membership and does not
assume $S_0\subseteq S$.

For either gap labeling $c$, expand the squared-gap sum as
$\sum_{a,b:\,c(a)=c(b)}w_aw_b$. The preceding implication preserves
every old pair whose two edges are uncovered. Consequently the decrease
in the squared-gap sum is at most the mass of all ordered pairs touching
$U$, namely
\[
 L^2-(L-C)^2=2LC-C^2.
\]
Divide by $2L^2$ and use $0\le C/L\le\min\{E/L,1\}$ together with
the monotonicity of $\psi$ on $[0,1]$. Finally $\psi(e)\le e\le E/L$
gives the original linear bound directly. The shared proof uses finite
partition sums and the previously verified actual circular gaps; the
real-coordinate adapter follows by telescoping edge lengths.
\end{outline}

This verifies approximation by anchored arcs while retaining the complete
source set in recovery. The quadratic saving increases the available
finite geometric credit; it does not establish economical approximations
for arbitrary tape geometries.

The anchoring premise matters: adding the isolated point $L/2$ to
$\{0\}$ lowers gap area from $1/2$ to $1/4$ despite adding zero
interval length. Small total measure without anchoring is insufficient.

For a phase $\rho$ and the first output time $s_j$ of column $j$, call
$(S_0,(I_i),E)$ an \emph{admissible cover witness} if
$\varnothing\ne S_0\subseteq\{0,\ldots,L-1\}$,
$\kappa_\rho(S_0)<\infty$, all arcs are anchored in $S_0$, their union
has length at most $E$, and
$S_\rho(s_j)\subseteq S_0\cup\bigcup_i I_i$.
Each formal witness supplies a common real subdivision containing the
arc endpoints and all raw coordinates, with the latter retaining their
values under an order-preserving embedding.
Choose any family of such witnesses and set
\begin{equation}\label{eq:cover-credit}
 H_{\rho,j}=\sup\left(\{0\}\cup
 \left\{\frac{(1-\alpha)_+^2}{2(2K+1)\kappa_\rho(S_0)}
           -\psi\!\left(\min\!\left\{\frac EL,1\right\}\right):
              \text{chosen admissible witnesses}\right\}\right).
\end{equation}
An empty witness family gives zero. The preceding lemmas show that
$0\le H_{\rho,j}\le\gaparea(S_\rho(s_j))\le1/2$, whether or not the
supremum is attained. Bounds below that subtract $E/L$ remain valid
relaxations of this stronger credit. In particular one may use the exact witness
$S_0=S_\rho(s_j)$, $E=0$. Let $\overline H$ be the arithmetic mean
of $H_{\rho,j}$ over the same legal phases and all $d$ columns.

\begin{theorem}[Theorem C: complete covers with the full paid cost]\label{thm:paid-source-cover}
\formal{SourceCovers/PaidSourceCover.lean}{SourceCovers.complete_cover_paid_deficit,SourceCovers.paid_deficit_of_gap_credit,SourceCovers.phaseCoverCredit_bounds,SourceCovers.meanCoverCredit}
Suppose every actual family outside the selected core and the $q$
exceptions has at most $a$ hops. For every integer $w\ge1$, define
\begin{align}\label{eq:full-paid-error}
 \mathsf E_q(w;J,M,R,a,\tau)
 ={}&\frac{2C_q(T-1)}{dL}+\frac{2(w-1)}L+\frac{W-b+1}{2L}\notag\\
 &+\frac\beta L\left[J+q+\frac\tau{M-R}
                    +\frac{(1+Na)C_q(T-1)}w\right].
\end{align}
Then the correctness deficit on the fixed full geometry satisfies
\begin{equation}\label{eq:paid-source-cover}
 \frac\Delta L\ge\overline H-\mathsf E_q(w;J,M,R,a,\tau).
\end{equation}
All payments in this inequality use the complete bank, with the
protected rank phase and the residual grid shift chosen compatibly.
The same bound holds with $\tau$ replaced by the hop mass of the
unprotected selected families that are nonleaves of the original core.
The finite statement does not require $T\le cL$.
\end{theorem}
\begin{outline}
At every step the formal parent relation contains the state predecessor
and every actual long work-tape hop; length-one work hops are state edges.
The selected relation uses the original core and exceptions, before phase
deletion. Pay the distinct outer lowers of those families and all actual
lowers deleted by the phase. Across every unpaid selected edge, exact-bank
source sets are nested. Every remaining live lower crosses its boundary
through a state edge or a hop of an unselected family. Full-span coloring
and the outside-family size bound therefore give at most $1+Na$ such lowers.

Index the exact list bank by its slots. This preserves both its raw value
sets and its summed variation, including coincident slots. For each legal
phase, \cref{lem:residual-source-recovery} supplies one shift whose residual
payment is at most $(1+Na)V_\rho/w$, where $V_\rho$ is that bank's actual
raw variation. The whole-record tuple at all distinct paid times is one
common transcript for every output. Local determinism and actual output
correctness imply functional recovery from this tuple and the complete raw
blocks. Thus \cref{thm:cyclic-freezing} applies to the same bank and tuple;
no global recovery or final payment estimate is assumed.

For each cover witness, telescoping in its supplied real subdivision
identifies weighted gap area with the original raw point-gap area. Apply
\cref{lem:fractional-gap,lem:anchored-erosion} to bound the witness score by
the complete frozen set's gap area. The scores, together with zero, are
bounded above by $1/2$, so their supremum satisfies the same bound without
an attainment premise. Neither the witness nor its credit depends on the
residual grid shift.

Replace the frozen areas by these credits and use
$V_\rho\le C_q(T-1)$. Average each phase's full signed inequality, keeping
credit and payment together. The mean initial payment is at most the
number of distinct selected outer lowers plus the unprotected nonleaf hop
mass divided by the actual number of legal phases. The outer count is at
most $J+q$, the hop mass at most $\tau$, and the phase count at least $M-R$.
Divide by $L$ to obtain
\cref{eq:full-paid-error,eq:paid-source-cover}. The formal proof retains the
actual variation and payment quantities before these relaxations.
\end{outline}

This theorem separates the two obligations of an impossibility proof:
establish a large geometric credit and keep the entire revelation bill
below it. Optimizing a witness is an analytic operation on the fixed
geometry, not a claim that the machine can compute it cheaply.

\begin{corollary}[Counting terminal core orbits]\label{cor:orbit-count-deficit}
\formal{SourceCovers/OrbitCountDeficit.lean}{SourceCovers.phaseCoreCover,SourceCovers.meanReciprocalCover,SourceCovers.meanReciprocalCover_bounds,SourceCovers.reciprocal_cover_paid_deficit,SourceCovers.orbit_count_paid_deficit}
Under the same hypotheses, put
\[
 \kappa_{\rho,j}=\kappa_\rho(S_\rho(s_j)),\qquad
 \mathcal R=\frac1{|\mathcal A|d}
       \sum_{\rho\in\mathcal A}\sum_{j=0}^{d-1}\frac1{\kappa_{\rho,j}},
\]
where $\mathcal A$ is the set of legal phases and $s_j$ is the first
output time of column $j$. Then
\begin{gather}
 \frac1{(2M+4)^{|\mathcal E|}}\le\mathcal R\le1,\notag\\
 \frac\Delta L\ge
 \frac{(1-\alpha)_+^2}{2(2K+1)}\mathcal R
 -\mathsf E_q(w;J,M,R,a,\tau).\label{eq:reciprocal-cover-deficit}
\end{gather}
In particular, since $|\mathcal E|\le q$,
\begin{equation}\label{eq:orbit-count-deficit}
 \frac\Delta L\ge
 \frac{(1-\alpha)_+^2}{2(2K+1)(2M+4)^q}
 -\mathsf E_q(w;J,M,R,a,\tau).
\end{equation}
\end{corollary}
\begin{outline}
The exact bank at time $s_j$ is a union of at most
$(\ell+1)^{|\mathcal E|}$ complete core orbits. Each covering root
belongs to its own orbit and hence to the bank. Regression puts it
at a time at most $s_j<T$, so it is a permitted root in the fractional
cover definition. Unit weights give
$1\le\kappa_{\rho,j}\le(2M+4)^{|\mathcal E|}<\infty$;
the lower bound uses the nonempty complete bank. Taking reciprocals
and averaging proves the bounds on $\mathcal R$.

Apply \cref{lem:fractional-gap} to each complete raw bank. Its credit
is $(1-\alpha)_+^2/[2(2K+1)\kappa_{\rho,j}]$, the score of the exact
zero-arc witness. The formal proof reuses the paid averaging argument
of \cref{thm:paid-source-cover}, with these gap lower bounds as its
credits. Factoring out the common numerator gives
\cref{eq:reciprocal-cover-deficit}; the lower bound on $\mathcal R$
then gives \cref{eq:orbit-count-deficit}.
\end{outline}

This completes the finite input to the necessary hop-tail theorem below.
Retaining the mean reciprocal can improve the finite deficit when actual
covers are cheaper than the terminal-orbit count. It preserves variation
across columns and phases without assuming economical covers for arbitrary
geometries or enlarging the uniform impossibility classes by itself.

\subsection{Two-exception cover suppliers}

Let $g$ and $e$ be respectively the newer and older exceptional maps,
ordered by decreasing owner; tied owners may be ordered either way.
If there is just one exception, take it as $e$ and let $g$ be the
identity; with no exceptions both maps are the identity.
Write
\begin{equation}\label{eq:older-closure}
 R_{e,\rho}(t)=A_\rho(t)\cup
                 \bigcup_{u\in A_\rho(t)}A_\rho(e(u)).
\end{equation}
This is the complete older-exception reachability relation. It is
reflexive, transitive and causal. With
$Z_\rho(t)=\{g(u):u\in A_\rho(t),\ g(u)\ne u\}$,
\begin{equation}\label{eq:two-exception-decomposition}
 O_\rho(t)=R_{e,\rho}(t)\cup
                       \bigcup_{z\in Z_\rho(t)}R_{e,\rho}(z).
\end{equation}
Indeed successful owners decrease strictly, so a path can use $g$
before $e$ but never conversely. Let $F_\rho(t)$ be the maximal elements
of $\{t\}\cup Z_\rho(t)$ under the order
$u\succeq_e v$ if $v\in R_{e,\rho}(u)$.
This is a partial order: a nonidentity step strictly decreases time.
The point $t$ is maximal and belongs to $F_\rho(t)$.

\begin{proposition}[Complete continuation suppliers]\label{prop:continuation-covers}
\notformalized
For two exceptions the following bounds hold for every phase and root.
\begin{enumerate}[label=(\alph*)]
\item $\kappa_\rho(S_\rho(t))\le(\ell+1)|F_\rho(t)|$.
\item Let $C_g(t)=A_\rho(t)\cup\bigcup_{u\in A_\rho(t)}A_\rho(g(u))$.
Remove from $\{e(v):v\in C_g(t),\ e(v)\ne v\}$ the points already in
$C_g(t)$, and let $s(t)$ be the number of maximal remaining points under
core reachability. Then $\kappa_\rho(S_\rho(t))\le\ell+1+s(t)$.
\item Partition $F_\rho(t)\setminus\{t\}$ into $r$ nonempty clusters
with representatives $a_1,\ldots,a_r$ and closed time hulls $I_1,\ldots,I_r$.
Fix a finite bank of time maps $(b_i)_i$ satisfying
$\{b_i(u):i\}=R_{e,\rho}(u)$ at every time $u$, as supplied by
\cref{thm:exact-source-bank}. These geometry-dependent maps are fixed
before any residual width or shift is chosen. Put
$E(t)=\sum_{v=1}^r\sum_i\Var(x\circ b_i;I_v)$.
An empty remainder has $r=E(t)=0$. Then
\begin{equation}\label{eq:cluster-gap}
 \gaparea(S_\rho(t))\ge
 \frac{(1-\alpha)_+^2}{2(2K+1)(\ell+1)(r+1)}-\frac{E(t)}L.
\end{equation}
\end{enumerate}
\end{proposition}
\begin{informalproof}
Every element of a finite partially ordered set is dominated by a maximal
one. Transitivity in \cref{eq:two-exception-decomposition} therefore replaces
its seed set by $F_\rho(t)$ without losing any endpoint. Each older closure
uses at most $\ell+1$ complete core orbits, proving (a).
For (b), first complete the core and newer exception. Any additional
endpoint is in a core orbit of an older landing seed. Core reachability
is transitive, and $C_g(t)$ is core-closed; seeds already inside it add
nothing. Maximal remaining seeds suffice, each contributing one core
orbit. The newer-only part contributes at most $\ell+1$.

For (c) use the base raw set
$S_0=x(R_{e,\rho}(t))\cup\bigcup_v x(R_{e,\rho}(a_v))$.
It has a cover by $(r+1)(\ell+1)$ complete core orbits. For each cluster
and older-bank slot, the raw range interval over $I_v$ contains its
representative value and has length at most the corresponding variation.
These intervals cover every raw source omitted from $S_0$, by (a)'s
exact closure identity. They are anchored in $S_0$ and their summed
length is at most $E(t)$. Apply
\cref{lem:fractional-gap,lem:anchored-erosion}.
\end{informalproof}

All three statements concern complete continuations after a return.
Replacing $x\circ b_i$ by the immediate input coordinate $x$ would remove
unaccounted older sources. These suppliers offer concrete sufficient
conditions for \cref{thm:paid-source-cover}; no automatic bound on their
values for all native geometries is asserted.
